% Options for packages loaded elsewhere
\PassOptionsToPackage{unicode}{hyperref}
\PassOptionsToPackage{hyphens}{url}
\documentclass[
  10pt,
  letterpaper,
]{article}
\usepackage{xcolor}
\usepackage[margin=23mm]{geometry}
\usepackage{amsmath,amssymb}
\setcounter{secnumdepth}{-\maxdimen} % remove section numbering
\usepackage{iftex}
\ifPDFTeX
  \usepackage[T1]{fontenc}
  \usepackage[utf8]{inputenc}
  \usepackage{textcomp} % provide euro and other symbols
\else % if luatex or xetex
  \usepackage{unicode-math} % this also loads fontspec
  \defaultfontfeatures{Scale=MatchLowercase}
  \defaultfontfeatures[\rmfamily]{Ligatures=TeX,Scale=1}
\fi
\usepackage{lmodern}
\ifPDFTeX\else
  % xetex/luatex font selection
\fi
% Use upquote if available, for straight quotes in verbatim environments
\IfFileExists{upquote.sty}{\usepackage{upquote}}{}
\IfFileExists{microtype.sty}{% use microtype if available
  \usepackage[]{microtype}
  \UseMicrotypeSet[protrusion]{basicmath} % disable protrusion for tt fonts
}{}
\makeatletter
\@ifundefined{KOMAClassName}{% if non-KOMA class
  \IfFileExists{parskip.sty}{%
    \usepackage{parskip}
  }{% else
    \setlength{\parindent}{0pt}
    \setlength{\parskip}{6pt plus 2pt minus 1pt}}
}{% if KOMA class
  \KOMAoptions{parskip=half}}
\makeatother
\usepackage{longtable,booktabs,array}
\newcounter{none} % for unnumbered tables
\usepackage{calc} % for calculating minipage widths
% Correct order of tables after \paragraph or \subparagraph
\usepackage{etoolbox}
\makeatletter
\patchcmd\longtable{\par}{\if@noskipsec\mbox{}\fi\par}{}{}
\makeatother
% Allow footnotes in longtable head/foot
\IfFileExists{footnotehyper.sty}{\usepackage{footnotehyper}}{\usepackage{footnote}}
\makesavenoteenv{longtable}
\setlength{\emergencystretch}{3em} % prevent overfull lines
\providecommand{\tightlist}{%
  \setlength{\itemsep}{0pt}\setlength{\parskip}{0pt}}
\usepackage[]{natbib}
\bibliographystyle{plainnat}
\usepackage{bookmark}
\IfFileExists{xurl.sty}{\usepackage{xurl}}{} % add URL line breaks if available
\urlstyle{same}
\hypersetup{
  pdftitle={Reviewer Response: Local Smoothing and Extrapolation},
  pdfauthor={Super Writer worked example},
  hidelinks,
  pdfcreator={LaTeX via pandoc}}

\title{Reviewer Response: Local Smoothing and Extrapolation}
\author{Super Writer worked example}
\date{September 2026}

\begin{document}
\maketitle

\textbf{Educational mock response.} These comments were written for the
repository; they did not come from a conference. The response uses only
the measurements in the companion toy regression example. This generic
layout is not a venue's official rebuttal template.

\subsection{R1: The robustness claim is broader than the
evidence}\label{r1-the-robustness-claim-is-broader-than-the-evidence}

\textbf{Comment.} The draft claims robustness to distribution shift,
although the extrapolation results appear worse for five neighbors.
Please reconcile the abstract with the table.

\textbf{Response.} We agree. The in-domain mean MSE decreases from
0.1083 to 0.0851, but the extrapolation mean increases from 2.6891 to
3.7012. We removed the general robustness claim. The revised abstract
reports both domains, and the Results section states that the in-domain
advantage does not extend to this tested extrapolation setting. The
Interpretation section explains why nearest-neighbor predictions are
constant beyond the rightmost training input, without claiming that this
determines the ranking on every function.

\subsection{R2: The uncertainty notation is
unclear}\label{r2-the-uncertainty-notation-is-unclear}

\textbf{Comment.} Are the error bars confidence intervals? Were grid
points treated as independent experimental runs?

\textbf{Response.} The table reports sample standard deviation across
five training seeds, not a confidence interval. Each seed produces one
MSE per method and domain by averaging over 101 fixed grid points. The
grid points are not counted as separate training replications. We now
state the aggregation unit in the Experimental Protocol and label the
column ``Sample SD.'' We also removed ``significantly'' from the
original claim because no significance test was run.

\subsection{R3: Add a parametric
baseline}\label{r3-add-a-parametric-baseline}

\textbf{Comment.} A parametric model could be a stronger extrapolation
baseline. Can the paper establish superiority over that alternative?

\textbf{Response.} The supplied experiment does not include that
comparison, so it cannot establish such superiority. We identify the
missing baseline explicitly in Limitations and retain the narrower
comparison between one and five neighbors. We have not added a result,
promised a completed run, or represented the suggested experiment as
evidence already available.

\subsection{Revision Locations}\label{revision-locations}

{\def\LTcaptype{none} % do not increment counter
\begin{longtable}[]{@{}
  >{\raggedright\arraybackslash}p{(\linewidth - 4\tabcolsep) * \real{0.3333}}
  >{\raggedright\arraybackslash}p{(\linewidth - 4\tabcolsep) * \real{0.3333}}
  >{\raggedright\arraybackslash}p{(\linewidth - 4\tabcolsep) * \real{0.3333}}@{}}
\toprule\noalign{}
\begin{minipage}[b]{\linewidth}\raggedright
Comment
\end{minipage} & \begin{minipage}[b]{\linewidth}\raggedright
Change
\end{minipage} & \begin{minipage}[b]{\linewidth}\raggedright
Location in the companion manuscript
\end{minipage} \\
\midrule\noalign{}
\endhead
\bottomrule\noalign{}
\endlastfoot
R1 & Report both directions; narrow robustness claim & Abstract;
Results; Interpretation \\
R2 & Identify seed-level sample SD; remove significance claim &
Experimental Protocol; Results table \\
R3 & State unevaluated baseline and limit comparison & Limitations \\
\end{longtable}
}

\end{document}
